% Part: model-theory
% Chapter: interpolation
% Section: interpolation-proof

\documentclass[../../../include/open-logic-section]{subfiles}

\begin{document}

\olfileid[pt]{mod}{int}{prf}

\olsection{Teorema da Interpolação de Craig}

\begin{thm}[Teorema da Interpolação de Craig]
\ollabel{thm:interpol} 
Se $\Entails !A \lif !B$, então existe !!a{sentence} $!C$ tal que
$\Entails !A \lif !C$ e $\Entails !C \lif !B$, e toda !!{constant},
!!{function} e !!{predicate} (exceto $\eq$) em $!C$ ocorre tanto em
$!A$ quanto em~$!B$. A !!{sentence} $!C$ é chamada de
\emph{interpolante} de $!A$ e~$!B$.
\end{thm}

\begin{proof}
Suponha que $\Lang{L}_1$ seja a linguagem de $!A$ e $\Lang{L}_2$, a
linguagem de $!B$. Seja
$\Lang{L}_0 = \Lang{L}_1 \cap \Lang{L}_2$. Para cada
$i \in \{0, 1, 2 \}$, obtenha $\Lang{L}'_i$ de $\Lang{L}_i$
acrescentando as infinitas novas !!{constant}s $\Obj c_0, \Obj c_1,
\Obj c_2, \dots$.

Se $!A$ é insatisfatível, $\lexists[x][\eq/[x][x]]$ é um
interpolante. Se $\lnot !B$ é insatisfatível (e, portanto, $!B$ é
válida), $\lexists[x][\eq[x][x]]$ é um interpolante. Assim, também
podemos supor que tanto $!A$ quanto $\lnot !B$ sejam satisfatíveis.

Para demonstrar a contrapositiva do Teorema da Interpolação, suponha
que não haja interpolante para $!A$ e $!B$. Em outras palavras,
suponha que $\{!A\}$ e $\{\lnot !B\}$ sejam inseparáveis em
$\Lang{L}_0$.

Nosso objetivo é estender o par $(\{ !A \}, \{\lnot!B\})$ a um par
maximalmente inseparável $(\Gamma^*, \Delta^*)$. Sejam $!A_0$, $!A_1$,
$!A_2$, \dots uma enumeração das !!{sentence}s de $\Lang{L}_1$, e
$!B_0$, $!B_1$, $!B_2$, \dots uma enumeração das !!{sentence}s
de~$\Lang{L}_2$. Definimos duas sequências crescentes de conjuntos de
!!{sentence}s $(\Gamma_n, \Delta_n)$, para $n \ge 0$, como segue.
Tome $\Gamma_0 = \{ !A\}$ e $\Delta_0 = \{\lnot !B \}$. Supondo que
$(\Gamma_n, \Delta_n)$ já estejam definidos, defina $\Gamma_{n+1}$ e
$\Delta_{n+1}$ por:
\begin{enumerate}
\item Se $\Gamma_n \cup \{!A_n \}$ e $\Delta_n$ são inseparáveis em
  $\Lang{L}'_0$, coloque $!A_n$ em $\Gamma_{n+1}$. Além disso, se
  $!A_n$ é uma !!{formula} existencial $\lexists[x][!S]$, escolha uma
  nova !!{constant} $c$ que não ocorra em $\Gamma_n$, $\Delta_n$,
  $!A_n$ nem em $!B_n$, e coloque $\Subst{!S}{c}{x}$ em
  $\Gamma_{n+1}$.
\item Se $\Gamma_{n+1}$ e $\Delta_n \cup \{!B_n \}$ são inseparáveis
  em $\Lang{L}'_0$, coloque $!B_n$ em $\Delta_{n+1}$. Além disso, se
  $!B_n$ é uma !!{formula} existencial $\lexists[x][!S]$, escolha uma
  nova !!{constant} $c$ que não ocorra em $\Gamma_{n+1}$, $\Delta_n$,
  $!A_n$ nem em $!B_n$, e coloque $\Subst{!S}{c}{x}$ em
  $\Delta_{n+1}$.
\end{enumerate}
Por fim, defina:
\begin{align*}
  \Gamma^* & = \bigcup_{n\ge 0} \Gamma_n, & 
  \Delta^* & = \bigcup_{n\ge 0} \Delta_n.
\end{align*}
Por indução simultânea em $n$, podemos agora demonstrar:
\begin{enumerate}
\item\ollabel{part-a} $\Gamma_n$ e $\Delta_n$ são inseparáveis em
  $\Lang{L}'_0$;
\item\ollabel{part-b} $\Gamma_{n+1}$ e $\Delta_n$ são inseparáveis em
    $\Lang{L}'_0$.
\end{enumerate}
A base de \olref{part-a} é dada por \olref[sep]{lem:sep1}. Para a parte
\olref{part-b}, precisamos distinguir três casos:
\begin{enumerate}
\item Se $\Gamma_0 \cup \{!A_0 \}$ e $\Delta_0$ são separáveis, então
  $\Gamma_1 = \Gamma_0$, e \olref{part-b} é simplesmente
  \olref{part-a};
\item Se $\Gamma_1 = \Gamma_0 \cup\{ !A_0\}$, então $\Gamma_1$ e
  $\Delta_0$ são inseparáveis por construção.
\item Resta considerar o caso em que $!A_0$ é existencial, de modo que
  $\Gamma_1 = \Gamma_0 \cup \{ \lexists[x][!S], \Subst{!S}{c}{x}
  \}$. Por construção, $\Gamma_0 \cup \{ \lexists[x][!S]\}$ e
  $\Delta_0$ são inseparáveis; por \olref[sep]{lem:sep2}, portanto,
  $\Gamma_0 \cup \{ \lexists[x][!S], \Subst{!S}{c}{x} \}$ e
  $\Delta_0$ também são inseparáveis.
\end{enumerate}
Isso completa a base da indução para \olref{part-a} e
\olref{part-b}. Passemos ao passo indutivo. Para \olref{part-a}, se
$\Delta_{n+1} = \Delta_n \cup \{ !B_n \}$, então $\Gamma_{n+1}$ e
$\Delta_{n+1}$ são inseparáveis por construção (mesmo quando $!B_n$ é
existencial, por \olref[sep]{lem:sep2}); se
$\Delta_{n+1} = \Delta_n$ (porque $\Gamma_{n+1}$ e
$\Delta_n \cup \{!B_n\}$ são separáveis), usamos a hipótese de indução
em \olref{part-b}. Para o passo indutivo de \olref{part-b}, se
$\Gamma_{n+2} = \Gamma_{n+1} \cup \{!A_{n+1} \}$, então
$\Gamma_{n+2}$ e $\Delta_{n+1}$ são inseparáveis por construção (mesmo
quando $!A_{n+1}$ é existencial, por \olref[sep]{lem:sep2}); se
$\Gamma_{n+2} = \Gamma_{n+1}$, usamos o caso indutivo de
\olref{part-a} que acabamos de demonstrar. Isso conclui a indução em
\olref{part-a} e \olref{part-b}.

Segue-se que $\Gamma^*$ e $\Delta^*$ são inseparáveis; caso contrário,
pela compacidade, existe $n \ge 0$ que separa $\Gamma_n$ e $\Delta_n$,
contra \olref{part-a}. Em particular, $\Gamma^*$ e $\Delta^*$ são
consistentes: pois, se o primeiro ou o segundo fosse inconsistente,
eles seriam separados, respectivamente, por
$\lexists[x][\eq/[x][x]]$ ou $\lforall[x][\eq[x][x]]$.

Mostremos agora que $\Gamma^*$ é maximalmente consistente em
$\Lang{L}'_1$ e, de modo análogo, $\Delta^*$ em $\Lang{L}'_2$. Para o
primeiro, suponha que $!A_n \notin \Gamma^*$ e
$\lnot !A_n \notin \Gamma^*$ para algum $n \ge 0$. Se
$!A_n \notin \Gamma^*$, então $\Gamma_n \cup \{!A_n \}$ é separável de
$\Delta_n$, e portanto existe $!C \in \Lang{L}'_0$ tal que:
\begin{align*}
  \Gamma^* & \Entails !A_n \lif !C, & 
  \Delta^* & \Entails \lnot !C.
\end{align*}
Do mesmo modo, se $\lnot !A_n \notin \Gamma^*$, existe
$!C' \in \Lang{L}'_0$ tal que:
\begin{align*}
  \Gamma^* & \Entails \lnot !A_n \lif !C', & 
  \Delta^* & \Entails \lnot !C'.
\end{align*}
Pela lógica proposicional, $\Gamma^* \Entails !C \lor !C'$ e
$\Delta^* \Entails \lnot (!C \lor !C')$; logo, $!C \lor !C'$ separa
$\Gamma^*$ e $\Delta^*$. Um argumento semelhante estabelece que
$\Delta^*$ é maximal.

Por fim, mostramos que $\Gamma^* \cap \Delta^*$ é maximalmente
consistente em $\Lang{L}'_0$. É evidentemente consistente, pois é a
interseção de conjuntos consistentes. Para mostrar a maximalidade, seja
$!S \in \Lang{L}'_0$. $\Gamma^*$ é maximal em
$\Lang{L'_1} \supseteq \Lang{L'_0}$, e, de modo análogo, $\Delta^*$ é
maximal em $\Lang{L'_2} \supseteq \Lang{L'_0}$. Segue-se que ou
$!S \in \Gamma^*$ ou $\lnot !S \in \Gamma^*$, e ou
$!S \in \Delta^*$ ou $\lnot !S \in \Delta^*$. Se $!S \in \Gamma^*$ e
$\lnot !S \in \Delta^*$, então $!S$ separaria $\Gamma^*$ e $\Delta^*$;
se $\lnot !S \in \Gamma^*$ e $!S \in \Delta^*$, então $\Gamma^*$ e
$\Delta^*$ seriam separados por $\lnot !S$. Portanto, ou
$!S \in \Gamma^* \cap \Delta^*$ ou
$\lnot !S \in \Gamma^* \cap \Delta^*$, e
$\Gamma^* \cap \Delta^*$ é maximal.

Como $\Gamma^*$ é maximalmente consistente, ele tem um modelo
$\Struct{M}'_1$ cujo !!{domain} $\Domain{M'_1}$ compreende todos e
somente os elementos $\Assign{c}{M'_1}$ que interpretam as
!!{constant}s---exatamente como na demonstração do teorema da
completude (\olref[fol][com][cth]{thm:completeness}). Do mesmo modo,
$\Delta^*$ tem um modelo $\Struct{M}'_2$ cujo !!{domain}
$\Domain{M'_2}$ é dado pelas interpretações $\Assign{c}{M'_2}$ das
!!{constant}s.

Obtenha $\Struct{M_1}$ de $\Struct{M'_1}$ descartando as interpretações
das !!{constant}s, !!{function}s e !!{predicate}s em
$\Lang{L'_1} \setminus \Lang{L'_0}$, e proceda de modo análogo para
$\Struct{M_2}$. Então a aplicação $h \colon M_1 \to M_2$ definida por
$h(\Assign{c}{M'_1}) = \Assign{c}{M'_2}$ é um isomorfismo em
$\Lang{L}'_0$, pois, como mostramos, $\Gamma^* \cap \Delta^*$ é
maximalmente consistente em $\Lang{L}'_0$. Isso se segue porque toda
$\Lang{L}'_0$-!!{sentence} ou pertence a ambos $\Gamma^*$ e
$\Delta^*$, ou não pertence a nenhum deles: assim,
$\Assign{c}{M'_1} \in \Assign{P}{M'_1}$ se, e somente se,
$\Atom{P}{c} \in \Gamma^*$ se, e somente se,
$\Atom{P}{c} \in \Delta^*$ se, e somente se,
$\Assign{c}{M'_2} \in \Assign{P}{M'_2}$. As demais condições
satisfeitas por isomorfismos podem ser estabelecidas de modo análogo.

Definamos agora um modelo $\Struct{M}$ para a !!{language}
$\Lang{L_1} \cup \Lang{L_2}$ como segue:
\begin{enumerate}
\item O !!{domain} $\Domain{M}$ é simplesmente $\Domain{M_2}$, isto
  é, o conjunto de todos os elementos $\Assign{c}{M'_2}$;
\item Se !!a{predicate}~$P$ pertence a
  $\Lang{L_2} \setminus \Lang{L_1}$, então
  $\Assign{P}{M} = \Assign{P}{M'_2}$;
\item Se um predicado $P$ pertence a
  $\Lang{L}_1\setminus \Lang{L}_2$, então
  $\Assign{P}{M} = h(\Assign{P}{M'_2})$, isto é,
  \[
  \tuple{\Assign{c_1}{M'_2}, \dots, \Assign{c_n}{M'_2}} \in
  \Assign{P}{M}
  \]
  se, e somente se,
  \[
  \tuple{\Assign{c_1}{M'_1}, \dots, \Assign{c_n}{M'_1}} \in
  \Assign{P}{M'_1}.
  \]
\item Se !!a{predicate} $P$ pertence a $\Lang{L}_0$, então $\Assign{P}{M} =
  \Assign{P}{M'_2} = h(\Assign{P}{M'_1})$. 
\item As !!^{function}s de $\Lang{L}_1 \cup \Lang{L}_2$, inclusive as
  !!{constant}s, são tratadas de modo análogo.
\end{enumerate}

Por fim, mostra-se por indução nas !!{formula}s que $\Struct{M}$
concorda com $\Struct{M'_1}$ em todas as !!{formula}s de $\Lang{L'_1}$
e com $\Struct{M'_2}$ em todas as !!{formula}s de $\Lang{L'_2}$. Em
particular, $\Struct{M} \Entails \Gamma^* \cup \Delta^*$, donde
$\Struct{M} \Entails !A$, $\Struct{M} \Entails \lnot!B$ e
$\not\Entails !A \lif !B$. Isso conclui a demonstração do Teorema da
Interpolação de Craig.
\end{proof}

\end{document}
